\documentclass[11pt,a4paper]{article}
\usepackage[UTF8]{ctex}
\usepackage{geometry}
\geometry{left=2.5cm,right=2.5cm,top=2.5cm,bottom=2.5cm}
\usepackage{amsmath,amssymb,amsthm}
\usepackage{hyperref}
\usepackage{booktabs}
\usepackage{enumitem}
\usepackage{xltabular}
\hypersetup{colorlinks=true,linkcolor=blue,urlcolor=blue}

\newtheorem{definition}{定义}[section]
\newtheorem{proposition}{命题}[section]
\newtheorem{remark}{注记}[section]

\title{AI 化攻防：弹药质量维度深化（修正版）\\
——分层分配、Stackelberg 非对称性与柯克霍夫反向压制}
\author{李龙胜}
\date{\today}

\begin{document}

\maketitle

\begin{center}
\textit{
弹药质量不再是结语中的一句隐喻。\\
高A值告警获得高质量攻防投入，低A值告警获得低质量投入——\\
质量预算的分配规则与既有分析能力分配完全同构。\\
攻击方在质量维度拥有Stackelberg先发优势——\\
但防御方可以通过“快”而非“强”来夺回先手。\\
微调不是为了更强，而是为了更快。\\
柯克霍夫原则在AI时代的保护对象\\
从数值参数扩展为几何结构。
}
\end{center}

\vspace{5mm}

\begin{abstract}
本文是对AI化攻防专题卷中“弹药质量取代弹药数量”洞察的系统深化。
首先建立弹药质量的形式化模型——攻击生成质量与防御检测质量
在严格凸成本函数下形成唯一纳什均衡。
进而将质量预算在告警类别间按A值降序分配，
其最优分配结构与既有框架完全同构。
揭示质量维度的Stackelberg非对称性——
攻击方的生成质量直接压制防御方有效检测率，
但防御方通过提高微调频率夺回先发优势。
最后提出柯克霍夫反向压制策略：
防御方通过周期性微调模型特征空间，
将攻击方的先发优势反向瓦解。
本文经外部审查修正，补充了质量成本函数凸性的来源分析、
有效检测率交互的非线性讨论、Stackelberg先发优势的动态交替，
以及特征空间微调的工程可行性。
\end{abstract}

\tableofcontents

\section{引言：从隐喻到形式化}

在AI化攻防专题卷和批判文档中，
“弹药质量取代弹药数量”作为一个核心洞察被提出。
当攻防双方都部署大模型、推理成本趋近于零时，
攻击方的竞争力不再体现为“能生成多少攻击payload”，
而是“生成的payload有多难被检测”；
防御方的竞争力不再体现为“能分析多少条告警”，
而是“检测模型的精度有多高”。

本文将该洞察从定性描述展开为独立的形式化模型。
质量维度的成本结构、均衡条件、预算分配
与防御方的反制策略在此得到完整阐述。

\section{弹药质量的基础形式化}

\subsection{质量成本函数的严格凸性}

攻击方生成质量为 \(q_A \in [0,1]\) 的对抗样本——
\(q_A\) 度量该样本逃逸防御方检测的概率。
生成成本函数 \(C_A(q_A)\) 满足
\[
C_A(0) = 0, \quad C_A'(q_A) > 0, \quad C_A''(q_A) > 0,
\quad \lim_{q_A \to 1} C_A(q_A) = \infty.
\]
完美逃逸的生成成本趋于无穷。

防御方训练检测质量为 \(q_D \in [0,1]\) 的模型——
\(q_D\) 度量模型在正常样本上的检测率。
训练成本函数 \(C_D(q_D)\) 同样严格凸，
且 \(C_D(1) \to \infty\)——完美检测的训练成本趋于无穷。

\textbf{凸性的来源分析}：
防御方的训练成本随精度超线性增长，
在深度学习中有大量实证支持——
更大规模的数据、更长的训练时间、更频繁的对抗训练，
成本随精度提升而指数级增长。

攻击方的凸性则依赖于防御方模型的访问难度。
在黑盒场景下（防御方模型权重保密），
攻击方需通过查询来探测决策边界，
达到精度所需的查询次数随精度超线性增长，
单次查询消耗算力和封禁风险，
有效生成成本严格凸。
在白盒场景下（攻击方获取防御方模型权重），
攻击方可直接在本地计算梯度，
生成成本可能近似线性，凸性假设可能被削弱。

这意味着质量成本函数的凸性依赖于柯克霍夫原则在AI场景中的具体实施——
防御方必须保密模型权重，使攻击方始终处于黑盒状态。
模型权重成为新的秘密参数，
与第2卷的核心结论一致：
参数保密不仅保护防御方的裁定阈值，
在AI时代还保护质量成本函数的凸性结构本身。

\subsection{有效检测率衰减——交互项的非线性特征}

防御方的有效检测率受攻击方生成质量的衰减：
\[
q_D^{\text{eff}} = q_D \cdot (1 - \lambda q_A),
\]
其中 \(\lambda \in [0,1]\) 是攻击方质量对防御方检测率的衰减系数。

\textbf{线性衰减假设的边界}：
当前线性衰减对所有 \(q_D\) 水平施加相同的衰减比例。
在真实对抗中，攻击方质量和防御方质量的交互可能存在非线性——
当双方质量都处于中高水平时，交互最激烈；
当一方压倒另一方时，交互趋于平淡。
如果 \(q_A\) 很高但 \(q_D\) 很低，
防御方本来就检测不到，攻击方的质量优势是冗余的。
如果 \(q_A\) 很低但 \(q_D\) 很高，
防御方的检测精度几乎不受影响。
最激烈的攻防发生在双方质量都处于中高水平的区间。

更复杂的交互函数形式将在后续工作中探索，
当前线性衰减是双方质量交互的一阶近似。

这意味着AI攻防最激烈的战场在对等质量区间——
双方都投入了高质量AI，
没有一方能够轻易压倒另一方。
攻防双方的最优质量选择
可能因此形成错位均衡——
避免恰好与对方质量处于同一中高区间，
因为那里是消耗最大、收益最不确定的泥潭。

\subsection{双方收益函数与纳什均衡}

攻击方的收益为攻击成功的期望收益减去生成成本：
\[
U_A(q_A, q_D) = q_A \cdot (1 - q_D^{\text{eff}}) \cdot A_{\text{target}} - C_A(q_A).
\]
防御方的收益为负的漏报损失减去训练成本：
\[
U_D(q_A, q_D) = - (1 - q_D^{\text{eff}}) \cdot q_A \cdot A_{\text{target}} - C_D(q_D).
\]

双方在质量维度上形成纳什均衡 \((q_A^*, q_D^*)\)。
由双方成本函数的严格凸性，
一阶条件给出唯一解：
\[
C_A'(q_A^*) = (1 - q_D^* \cdot (1 - \lambda q_A^*)) \cdot A_{\text{target}},
\]
\[
C_D'(q_D^*) = q_A^* \cdot (1 - \lambda q_A^*) \cdot A_{\text{target}}.
\]
均衡存在且唯一。

\subsection{均衡的解析性质}

攻击方的均衡生成质量 \(q_A^*\) 随攻击目标价值 \(A_{\text{target}}\) 递增——
目标越值钱，攻击方投入的质量越高。
防御方的均衡检测质量 \(q_D^*\) 同样随 \(A_{\text{target}}\) 递增——
目标越值钱，防御方投入的检测质量也越高。

但防御方的均衡检测质量还受攻击方质量衰减系数 \(\lambda\) 的影响——
\(\lambda\) 越大，攻击方质量对防御方有效检测率的压制越强，
防御方在同一训练投入下能达到的有效检测率越低。
防御方的最优检测质量可能因此下降——
不是防御方不够努力，而是攻击方在质量维度上的先发优势
使防御方面临一个向上追赶的陡峭成本曲线。

\section{质量分层分配——高价值高精度，低价值低精度}

真实攻防中，质量是分层的。
不同告警类别的攻击价值 \(A_i\) 不同——
攻击方对高A值告警生成高质量对抗样本的边际收益更高，
防御方对高A值告警训练高质量检测模型的边际漏报损失避免也更高。

设攻击方的总质量生成预算为 \(B_A\)，
防御方的总质量训练预算为 \(B_D\)。
攻击方在预算约束下分配各告警类别的生成质量 \(\{q_{A,i}\}\)，
最大化总期望攻击收益：
\[
\max_{\{q_{A,i}\}} \sum_i q_{A,i} \cdot (1 - q_{D,i} \cdot (1 - \lambda q_{A,i})) \cdot A_i,
\quad \text{s.t.} \quad \sum_i C_A(q_{A,i}) \le B_A.
\]
防御方在预算约束下分配各告警类别的检测质量 \(\{q_{D,i}\}\)，
最小化总期望漏报损失：
\[
\min_{\{q_{D,i}\}} \sum_i (1 - q_{D,i} \cdot (1 - \lambda q_{A,i})) \cdot q_{A,i} \cdot A_i,
\quad \text{s.t.} \quad \sum_i C_D(q_{D,i}) \le B_D.
\]

最优分配均按 \(A_i\) 降序——
攻击目标价值越高的告警类别，分配越高的攻防质量。
这一结构与分析能力配额分配（按 \(\beta_i\) 降序）、
SOAR剧本优先级调度（按 \(V_j/c_j\) 降序）、
蜜罐部署密度分配（按 \(p_i/c_i\) 降序）
在数学上完全同构。
够用就行原则在质量维度上同样成立：
边际质量收益等于边际质量成本。

\section{质量 Stackelberg 非对称性与柯克霍夫反向压制}

\subsection{攻击方的先发优势与动态交替}

质量博弈存在一个根本的方向性不对称。
攻击方的生成质量 \(q_A\) 直接影响防御方的有效检测率
\(q_D^{\text{eff}} = q_D \cdot (1 - \lambda q_A)\)——
攻击方质量提升直接压制防御方的检测能力。
但防御方的检测质量 \(q_D\)
只能通过压缩攻击方的期望收益来间接影响攻击方的质量选择——
攻击方在下一轮调整 \(q_A\) 时才受到防御方质量提升的影响。

在单轮博弈中，攻击方拥有Stackelberg先发优势。

但在持续对抗训练的动态博弈中，
双方交替行动——
防御方每轮对抗训练后成为新一轮的先动者：
先发布新模型权重，攻击方再观测并生成对抗样本。
先发优势在双方之间周期性地交替持有。

这意味着柯克霍夫反向压制的核心
不是“微调幅度足够大”，而是“微调频率足够快”——
在攻击方完成新一轮质量校准之前，
防御方已经切换到下一个特征空间。
防御方不需要每一个回合都彻底改变特征空间，
只需要保证特征空间的切换频率
高于攻击方的质量校准速度。

防御方的最优微调频率
是使其在大多数回合中保持先发优势的那个频率。
这是够用就行原则在AI攻防节奏上的精确表达——
微调不是为了更强，而是为了更快。

\subsection{柯克霍夫反向压制——用不对称性对抗不对称性}

防御方通过周期性微调模型的特征空间和决策边界
来逆转攻击方的先发优势。

如果防御方在每个周期改变检测模型的特征空间——
通过对抗训练引入新的特征、随机化特征排序、
周期性替换模型架构或优化器初始化——
那么攻击方已生成的高质量对抗样本可能在特征空间的位移中失效。
攻击方不知道下一个周期的特征空间长什么样，
他在当前周期投入的高质量生成，在下个周期可能完全作废。

这是柯克霍夫原则在AI质量维度的直接应用。
在传统攻防中，柯克霍夫原则保护的核心是分析参数的不可预测性——
参数周期性扰动使攻击方的探测推断过期。
在AI时代，它保护的核心扩展为模型特征空间和决策边界的不可预测性——
周期性微调使攻击方的质量投入失效。
防御方不需要直接压制攻击方的质量优势，
而是用不对称性对抗不对称性——
攻击方质量先发优势，防御方特征空间不可预测性后发制人。
周期性微调不是直接提升 \(q_D\)，
而是增加攻击方达成有效逃逸的生成成本 \(C_A(q_A)\)——
攻击方不仅要生成高质量对抗样本，
还要在未知的特征空间中盲搜有效攻击方向。

\textbf{特征空间微调的工程可行性}：
彻底的模型架构替换实现高质量压制，
但成本极高且损害模型在正常样本上的性能。
轻量级扰动成本低但对高质量对抗样本效果有限。
最优微调幅度和频率需要在特征空间稳定性和不可预测性之间寻求均衡——
微调太少，攻击方适应；微调太多，防御方自身稳定性受损。
微调频率和幅度是两个耦合的决策变量：
频率高时幅度可以小——累积的小幅度变化让攻击方无法收敛；
频率低时幅度必须大——需要一次性制造足够大的位移来废弃攻击方的质量投资。
最优参数是边际压制收益等于边际稳定性损失的均衡点——
这是够用就行原则在AI防御中的又一次应用。

\section{诚实标注}

\begin{center}
\begin{xltabular}{\textwidth}{c|X|c}
\toprule
\textbf{命题} & \textbf{状态} & \textbf{层级} \\
\midrule
弹药质量成本函数的严格凸性 &
防御方训练成本凸性有大量实证支持；
攻击方凸性依赖于防御方模型保密性——
黑盒查询成本指数增长，白盒生成成本可能线性 &
II（白盒场景下凸性假设的削弱需实测数据验证） \\
有效检测率衰减的线性假设 &
线性衰减是双方质量交互的一阶近似；
更复杂的交互函数形式将在后续工作中探索 &
II（非线性交互的最激烈区域在双方质量相当时） \\
质量维度纳什均衡的存在性和唯一性 &
由双方成本函数严格凸性保证，一阶条件给出唯一解 &
II \\
质量分层分配的最优分配规则 &
按 \(A_i\) 降序分配质量预算，与已有框架完全同构 &
II \\
攻击方质量维度的先发优势 &
\(q_A\) 直接压制 \(q_D^{\text{eff}}\) 的机制已明确 &
II \\
Stackelberg先发优势的动态交替 &
在持续对抗训练中防御方每轮微调后成为先动者；
微调频率高于攻击方校准速度是夺回先手的关键 &
II（动态交替的多轮极限解待后续展开） \\
柯克霍夫反向压制——周期性微调特征空间 &
通过信息不对称增加攻击方质量生成成本；
最优微调频率和幅度是够用就行原则的又一次应用 &
II（特征空间微调的具体实现方案待工程验证） \\
弹药质量与弹药数量的联合优化 &
二维均衡框架已提出，联合优化的精确形式待后续展开 &
III \\
\bottomrule
\end{xltabular}
\end{center}

\section{结语}

弹药质量维度的深化揭示了四个层次。
分层分配——高价值告警获得高质量攻防投入，
质量预算按A值降序分配，与已有框架完全同构。
Stackelberg非对称性——攻击方在质量维度拥有先发优势，
但防御方通过提高微调频率夺回先手：
微调不是为了更强，而是为了更快。
柯克霍夫反向压制——防御方通过周期性微调特征空间
增加攻击方质量生成成本，逆转攻击方的先发优势。
柯克霍夫原则在AI时代的保护对象
从数值参数扩展为几何结构。

\vspace{5mm}
\noindent\textbf{文档信息}：
\begin{itemize}
    \item 作者：李龙胜
    \item 版本：修正版（整合外部审查反馈）
    \item 许可：CC BY 4.0
    \item 前置依赖：四卷体系、AI化攻防专题卷、
          AI化攻防批判修正版、攻防博弈论统一框架修正版
\end{itemize}

\end{document}